\section{Introduction}

Operational monitoring is intrinsically multivariate. Demand, embedded
generation, pumping, and frequency can each remain within a historically
familiar range while their joint configuration changes. Conversely, seasonal
or regime drift can make many conventional anomaly scores large without a
discrete operational event. The practical question is therefore not only
whether a value is extreme, but whether a window is unusual relative to the
information that would actually have been available when it arrived.

That question is difficult for three reasons. First, telemetry is serially
dependent and nonstationary, whereas simple rank arguments are sharpest under
exchangeability \citep{shafer2008tutorial}. Second, coarse summaries may erase
short physical transients even when they make longer multi-stream episodes
easy to rank. Third, an operational detector includes a selection policy as
well as a score. A nominal multiple-testing procedure, a fallback rule, and an
attention budget make different claims and must be reported separately.

TRACE-C retains its historical name, ``temporal relational anomaly detection
with copula calibration,'' but this manuscript narrows that description to the
shipped algorithm. The method performs strictly-prior online scoring with
rolling robust-z marginals and rank calibration. Its relational channel has the
algebraic form of a Gaussian copula log-density contrast \citep{song2000gaussiancopula},
but it is evaluated on robust-z residuals without a probability-integral or
normal-score transform. It is therefore not a literal copula density. We keep
the TRACE-C name for continuity and use ``rank-calibrated relational anomaly
detection'' as the scientifically accurate description.

The output $p$-values have an equally limited interpretation. They are
upper-tail conformal/rank values used to order and select windows. They are not
the probability that an alert is genuine, not posterior confidence, and not a
guarantee of time-series coverage. Exchangeability provides a useful benchmark
for interpreting ranks, but it is not established for these data.

This paper makes four contributions:
\begin{enumerate}
  \item It specifies an auditable three-channel detector that preserves
  residual magnitude while making every reference strictly prior to the
  scored window.
  \item It distinguishes the Gaussian copula-form dependence contrast actually
  implemented from literal copula calibration, and distinguishes the Fisher
  aggregation score from a chi-square test.
  \item It reports the full selection path---ordinary BH attempted first, a
  disclosed record-rule fallback, and a fixed-block alert budget---together
  with the assumptions that prevent a nominal FDR claim.
  \item It reports a frozen 2020 run, a post-hoc non-identical baseline
  comparison, and a 2019 development-year channel ablation showing that
  Atiyah's rank-1 result is local-led rather than copula-only, and that
  Fisher fusion can bury the short frequency event relative to the temporal
  channel alone.
\end{enumerate}

The results precede the limitations and conclusion so that the honesty ledger
can be read against concrete evidence. TRACE-C ranks Storm Atiyah first in
2019 development, but the ablation does not support a joint-surprise account
of that window. Reconstruction methods isolate the short 2019 frequency event
better. The frozen 2020 run produces no selected alerts. These outcomes define
the operating envelope supported by the present evidence.
